\section{What derivative moments would have to prove}
\label{sec:fd-moments}
This section turns derivative moments into an exact finite inequality and records the quantifiers needed for simplicity at every height.
For a finite multiset $\mathcal Z$ of nontrivial zeta zeros, with each zero
repeated according to its analytic multiplicity, put
\[
 u_\rho=|\zeta'(\rho)|^2,\quad
 s_k=\sum_{\rho\in\mathcal Z}u_\rho^k\ (k\ge1),\quad
 s_0=|\mathcal Z|.
\]
Let $M_0$ be the sum of the multiplicities of the multiple zeros in
$\mathcal Z$.  It is zero or an integer at least two.
Let $r$ be the number of distinct positive values among the $u_\rho$.
Define the Christoffel variational quantity
\[
 \lambda_d=\min_{P\in\mathbb R[X],\,\deg P\le d,\,P(0)=1}
                  \sum_{\rho\in\mathcal Z}P(u_\rho)^2.
\]
This is the classical Christoffel function at the origin for the finite
measure $\sum_\rho\delta_{u_\rho}$.
The variational and inverse-moment-matrix descriptions are classical
\cite{LasserrePauwels2019}; the same operation appears in the discussion
of moment-based inertia counts in Claude's simple-zero paper
\cite[Section~7.5(d),(f)]{Claude2026}.
The finite derivative measure here permits the following exact endpoint formula.

\begin{proposition}[A finite moment certificate for simplicity]
\label{prop:fd-moment}
For $0\le d\le r$, with the determinant of an empty matrix taken as one,
\[
 \lambda_d=
 \frac{\det[s_{i+j}]_{i,j=0}^{d}}
      {\det[s_{i+j}]_{i,j=1}^{d}},\qquad
 \lambda_r=M_0.
\]
The denominator is positive.  The sequence $\lambda_d$ is nonincreasing,
and $\lambda_d\ge M_0$ for all $d$.  The zeros in $\mathcal Z$ are
all simple precisely when one of these finite tests has $\lambda_d<2$.
At degree one,
\[
 \lambda_1=s_0-\frac{s_1^2}{s_2},
\]
whenever $r\ge1$; this is the usual Cauchy--Schwarz simple-zero count.
\end{proposition}

\begin{proof}
Write the measure as
$\mu=M_0\delta_0+\sum_{j=1}^r w_j\delta_{v_j}$,
where the $v_j>0$ are distinct and $w_j$ are positive integers.
Every admissible polynomial has squared integral at least $M_0$.
The polynomial $P(X)=\prod_{j=1}^r(1-X/v_j)$ attains this lower
bound and has degree $r$.  Enlarging the admissible polynomial space
makes the minima nonincreasing.
For $d\le r$, the lower moment block
$C=[s_{i+j}]_{i,j=1}^d$ is positive definite: its quadratic form is
$\sum_jw_j(\sum_{i=1}^d c_i v_j^i)^2$.
A nonzero polynomial $XQ(X)$ of degree at most $d$ cannot vanish at
$r\ge d$ distinct positive points in addition to zero.
For $P=1+\sum_{i=1}^d c_iX^i$, the objective is
$s_0+2b^Tc+c^TCc$, where $b=(s_1,\ldots,s_d)^T$.
Its minimum is $s_0-b^TC^{-1}b$.
The block determinant formula gives the displayed quotient, including
the case where the full moment matrix is singular.
At $d=0$ the formula reads $\lambda_0=s_0$.
When $r=0$ all the mass is at zero and this already equals $M_0$.
Finally a multiple analytic zero contributes at least two to $M_0$.
Thus $\lambda_d<2$ forces $M_0=0$, and $M_0=0$ gives
$\lambda_r=0$.  The empty multiset has $\lambda_0=0$ and also
satisfies the conclusion.
\end{proof}

The moments can also be evaluated on the same fixed contours used
for the heat minors.  Let $f(z)=H_0(z)$ and let a positively oriented
contour $\Gamma$ enclose a finite set of positive endpoint zeros,
with no zero on its boundary.  Alp\"oge's theorem makes those zeros
real, so $f'(x)^2\ge0$.  Put
\begin{equation}\label{eq:fd-contourmom}
 s_k^{\Gamma}=\frac1{2\pi i}\int_\Gamma
            f'(z)^{2k}\frac{f'(z)}{f(z)}\,dz,\qquad k\ge0.
\end{equation}
The residue at a zero $x$ of multiplicity $m_x$ is
$m_x f'(x)^{2k}$.  Hence~\eqref{eq:fd-contourmom} gives exactly
the moments of the positive finite measure
$\sum_xm_x\delta_{f'(x)^2}$.  The proposition applies unchanged,
and counts exactly the same multiple endpoints.
This is an explicit relation between the moment certificate and
contour integration, through the holomorphic observable $(f')^2$.
The two derivative normalizations are related at
$\rho=1/2+i\gamma$, $x=2\gamma$, by
\[
 f'(2\gamma)=\frac{i}{16}A(\rho)\zeta'(\rho),\qquad
 A(s)=\tfrac12s(s-1)\pi^{-s/2}\Gamma(s/2),
\]
so
\[
 f'(2\gamma)^2=|f'(2\gamma)|^2
             =\frac{|A(\rho)|^2}{256}|\zeta'(\rho)|^2.
\]
The factor is positive at every nontrivial zero.  Both measures
therefore have the same atom at zero, with its full multiplicity,
although the positive atoms have the displayed different locations.

A rigorous numerical certificate need not invert a nearly singular
moment matrix.  Choose a specific real polynomial
$P(X)=\sum_{i=0}^d c_iX^i$, with $c_0=1$.
Suppose approximate moments obey
$|s_k-\widetilde s_k|\le\epsilon_k$ for $0\le k\le2d$.
Then
\begin{equation}\label{eq:fd-momentcert}
 \sum_{\rho\in\mathcal Z}P(u_\rho)^2
 \le\sum_{i,j=0}^dc_ic_j\widetilde s_{i+j}
       +\sum_{i,j=0}^d|c_ic_j|\epsilon_{i+j}.
\end{equation}
This follows by expanding the square and applying the triangle
inequality to each moment error.  An upper bound less than two
certifies simplicity throughout the contour.
For the derivative values of $f$, every $s_k$ on the right is the
contour integral~\eqref{eq:fd-contourmom}.
Thus contour separation and bounds for $f,f'$ and their approximation
errors yield a finite certificate using the same boundary data as the
heat-minor method.  Its numerical threshold is an absolute number.

The exact all-height problem in this formulation is
\[
 \text{for every }T>0\text{ there are }d\ge0\text{ and }P(0)=1
 \text{ with }
 \sum_{0<\Im\rho\le T}P(|\zeta'(\rho)|^2)^2<2.
\]
A proved simple-zero proportion supplies an upper bound of order
$N(T)$ for the corresponding count of remaining multiplicities.
The displayed criterion requires a bound below two, at each height.
Equation~\eqref{eq:fd-momentcert} specifies the error precision that
would turn a moment estimate into that integer-valued certificate.
No estimate reaching this threshold at every height is proved here.

\section{Finite zero surgery and the information in leading moments}
\label{sec:fd-surgery}
This section gives an exact transformation preserving the leading positive-moment data while creating a prescribed multiple endpoint.
Choose two known simple positive ordinates $u<v$, for example the
first two, and set
\[
 R(s)=\frac{(s-1/2)^2+v^2}{(s-1/2)^2+u^2},\qquad
 \widetilde\zeta(s)=R(s)\zeta(s),\qquad
 \widetilde\xi(s)=R(s)\xi(s).
\]
The apparent poles at $1/2\pm iu$ are removable by simplicity.
The map removes those two zeros and adds one further copy of each
zero at $1/2\pm iv$.
For every $T>v$, the positive-ordinate zero counts, including
multiplicity, satisfy
\[
 \widetilde N(T)=N(T),\qquad
 \widetilde N_{\rm simple}(T)=N_{\rm simple}(T)-2.
\]
The completed function is entire, of the same order, and obeys
$\widetilde\xi(1-s)=\widetilde\xi(s)$ and the same conjugation law,
because $R(1-s)=R(s)$ and $R(\overline s)=\overline{R(s)}$.
Its nontrivial zeros all lie on the critical line by Alp\"oge's theorem
and the explicit zero replacement.  In the variable $z$ of $H_0$,
\[
 \frac{\widetilde H_0(i\sqrt w)}{\widetilde H_0(0)}
 =\frac{H_0(i\sqrt w)}{H_0(0)}
       \frac{1+w/(4v^2)}{1+w/(4u^2)}.
\]
The right side cancels one positive real product parameter and
repeats another; its endpoint Taylor coefficients remain positive.
This identity states exactly which coefficient object the
transformation produces.

For every other positive-ordinate zero,
\begin{equation}\label{eq:fd-surgeryder}
 \widetilde\zeta'(1/2+i\gamma)
  =\frac{\gamma^2-v^2}{\gamma^2-u^2}\,
                          \zeta'(1/2+i\gamma).
\end{equation}
This includes multiple zeros, where both derivatives vanish.
For $p>0$ put $M_p(T)=\sum_{0<\Im\rho\le T}|\zeta'(\rho)|^p$
and define $\widetilde M_p$ with the changed zero multiset.
For every $Y>2v$ and every $T\ge Y$,
\begin{equation}\label{eq:fd-momentsurgery}
 |\widetilde M_p(T)-M_p(T)|
 \le C_{p,Y}+\delta_p(Y)M_p(T),\qquad
 \delta_p(Y)=1-\left(1-\frac{v^2-u^2}{Y^2-u^2}\right)^p.
\end{equation}
Here $C_{p,Y}$ is the finite sum of absolute moment differences below
$Y$, including the two replaced endpoints.
To prove the inequality, for $\gamma\ge Y$ the positive ratio in
\eqref{eq:fd-surgeryder} lies between
$1-(v^2-u^2)/(Y^2-u^2)$ and one.
Raise it to the power $p$, multiply by $|\zeta'(\rho)|^p$, and
sum.  Since $\delta_p(Y)\to0$, every divergent leading positive
moment asymptotic is transported with relative error tending to zero.
The exact counting identities likewise transport every asymptotic
simple-zero proportion.  Nonetheless the zeros at $v$ are double.
These are explicit limitations of that collection of asymptotic data.

The transformation also identifies an arithmetic property that it changes.
As $\sigma\to+\infty$ along the real axis,
\[
 \widetilde\zeta(\sigma)
 =1+\frac{v^2-u^2}{\sigma^2}+O(\sigma^{-3})+O(2^{-\sigma}).
\]
It therefore has no ordinary Dirichlet series with leading coefficient
one that is absolutely convergent in a right half-plane.
Indeed, absolute convergence at $\sigma_0$ would imply
\[
 \left|\sum_{n\ge2}a_n n^{-\sigma}\right|
 \le2^{-(\sigma-\sigma_0)}
        \sum_{n\ge2}|a_n|n^{-\sigma_0}=O(2^{-\sigma}),
\]
contradicting the nonzero $\sigma^{-2}$ term.
Thus the exact transformation preserves the stated endpoint and
leading-moment information while changing the displayed Dirichlet-series
property.  A proof of simplicity using arithmetic information must use
more than the transported data.

\subsection{The positive kernel under finite zero surgery}
\label{sec:fd-positive-surgery}

The same transformation has an explicit positive kernel. We first
give the kernel estimate needed for its construction. The classical
properties of the Riemann kernel are collected in
\cite[Theorem~4.1, p.~10]{Csordas2015}; the monotonicity estimate is
proved here so that the positivity argument has a complete local input.

\begin{lemma}
\label{lem:fd-kernel-decreasing}
The Riemann kernel \(\Phi\) is strictly decreasing on \((0,\infty)\).
It is even and analytic in \(|\Im u|<\pi/8\), and every derivative
is \(O(\exp(-c e^{4|u|}))\) on the real axis for every \(0<c<\pi\).
\end{lemma}

\begin{proof}
For \(u\ge0\), write \(y_n=\pi n^2e^{4u}\). The defining series is
\[
 \Phi(u)=\sum_{n\ge1}e^u(2y_n^2-3y_n)e^{-y_n}.
\]
It and its derivatives converge locally uniformly in the strip,
since \(\Re(e^{4u})>0\) there. The stated decay follows by
absorbing the polynomial-exponential factors into
\(\exp((\pi-c)e^{4u})\). To verify evenness, let
\(\theta(t)=\sum_{n\in\mathbb Z}e^{-\pi n^2t}\).
The Gaussian Poisson identity gives
\(\theta(t)=t^{-1/2}\theta(1/t)\). Hence
\(\Theta(u)=e^u\theta(e^{4u})\) is even, and termwise
differentiation gives
\(\Phi=(\Theta''-\Theta)/16\). In particular \(\Phi'(0)=0\).

For the individual summands \(\Phi_n\), direct differentiation gives
\[
 \Phi_n'(u)=e^u y_ne^{-y_n}p(y_n),\qquad
 \Phi_n''(u)=e^u y_ne^{-y_n}Q(y_n),
\]
where
\[
 p(y)=-8y^2+30y-15,\qquad
 Q(y)=32y^3-224y^2+330y-75.
\]
The larger root of \(p\) is
\((15+\sqrt{105})/8<79/25\).
For \(u\ge1/500\), the elementary bounds
\(\pi>157/50\) and \(e^{1/125}>126/125\) give
\(y_n>79/25\). Every summand of \(\Phi'\) is then negative.

It remains to consider \(0\le u\le1/500\).
Here \(3<y_1<16/5\), using \(\pi<22/7\) and
\(e^{1/125}<125/124\).
On \([3,16/5]\),
\[
 Q'(y)=96y^2-448y+330
 \le96(16/5)^2-448\cdot3+330<0.
\]
Thus \(Q(y_1)\le Q(3)=-237\).
Since \(e<3\), we have \(y_1e^{-y_1}>1/27\), so
\(\Phi_1''<-79/9\).
For \(n\ge2\), one has \(y_n\ge12\),
\(y_n<(16/5)n^2\), \(e^u<101/100\), and
\(Q(y_n)<35y_n^3\). Therefore
\[
 \sum_{n\ge2}\Phi_n''(u)
 <\frac{101}{100}\,35(16/5)^4
                 \sum_{n\ge2}n^8e^{-3n^2}.
\]
Successive terms of the last sum have ratio at most
\((3/2)^8e^{-15}<1/10000\). Using \(e>8/3\), its first term
and geometric tail give
\[
 \sum_{n\ge2}n^8e^{-3n^2}
 <256(3/8)^{12}\frac{10000}{9999}<\frac1{500}.
\]
Finally
\(\frac{101}{100}35(16/5)^4/500<79/9\).
Thus \(\Phi''<0\) on this short interval.
Together with \(\Phi'(0)=0\), this proves strict decrease everywhere.
\end{proof}

\begin{proposition}[A positive kernel for the exact replacement]
\label{prop:fd-positive-surgery}
Let \(U=2u\) and \(V=2v\) be the two simple positive endpoints
used in the finite surgery, so \(0<U<V\). Define, for \(x\ge0\),
\[
 \Psi_U(x)=U\int_x^\infty\sin(U(t-x))\Phi(t)\,dt,
 \qquad
 \widetilde\Phi(x)=\Phi(x)+
                   \left(\frac{V^2}{U^2}-1\right)\Psi_U(x),
\]
and extend both functions evenly to the real axis.
Then \(\Psi_U\) and \(\widetilde\Phi\) are strictly positive,
even, smooth, and analytic in \(|\Im x|<\pi/8\).
Every real-axis derivative has superexponential decay of the
form in Lemma~\ref{lem:fd-kernel-decreasing}. Their cosine transforms
satisfy the exact identities
\[
 \int_0^\infty\Psi_U(x)\cos(zx)\,dx
       =\frac{H_0(z)}{1-z^2/U^2},\qquad
 \int_0^\infty\widetilde\Phi(x)\cos(zx)\,dx
       =\frac{V^2-z^2}{U^2-z^2}H_0(z).
\]
All apparent singularities are removable. In particular the entire
heat family
\[
 \widetilde H_t(z)=\int_0^\infty
                e^{t x^2}\widetilde\Phi(x)\cos(zx)\,dx
\]
is defined for every complex \(t,z\), and its endpoint is precisely
the completed-function surgery already described.
For real \(a\ge0\), every coefficient of
\(\widetilde H_{-a}(i\sqrt w)/\widetilde H_{-a}(0)\) is positive.
\end{proposition}

\begin{proof}
Put \(b=\pi/U\). Pairing successive half periods gives
\[
 \frac{\Psi_U(x)}U
 =\int_0^b\sin(Ut)\sum_{j\ge0}
  \bigl[\Phi(x+t+2jb)-\Phi(x+t+(2j+1)b)\bigr]dt>0.
\]
The kernel's strict decrease proves the sign; its decay justifies
the pairing and exchange of sum and integral.
Differentiating the defining integral gives
\[
 \Psi_U'(x)=-U^2\int_x^\infty\cos(U(t-x))\Phi(t)\,dt,
 \qquad
 \Psi_U''+U^2\Psi_U=U^2\Phi.
\]
At the origin, \(\Psi_U'(0)=-U^2H_0(U)=0\).
The differential equation with even analytic forcing therefore
has an even analytic solution in the same strip. This proves the
claimed regularity of the even extension.

On the real axis, \(|\sin(Uv)|\le Uv\) gives
\[
 |\Psi_U(x)|\le U^2\int_0^\infty v\Phi(x+v)\,dv.
\]
The defining theta series gives
\(\Phi(y)\le C e^{9y-\pi e^{4y}}\) for \(y\ge0\).
Since \(e^{4v}-1\ge4v\), it follows that
\[
 |\Psi_U(x)|\le
 \frac{CU^2 e^{9x-\pi e^{4x}}}{(4\pi e^{4x}-9)^2}
       =O_U(e^{x-\pi e^{4x}}).
\]
The first derivative has the analogous estimate from its integral;
the differential equation and induction give the bound for every
higher derivative. Thus integration by parts has no boundary
contribution at infinity, and its contribution at zero vanishes
because \(\Psi_U'(0)=0\).
If \(K_U\) is the cosine transform of \(\Psi_U\), the equation
therefore gives
\((U^2-z^2)K_U(z)=U^2H_0(z)\).
Both transforms are entire by the decay just established, proving
the first identity and removability at \(\pm U\).
Adding the two kernels gives the second identity because
\[
 1+\frac{V^2/U^2-1}{1-z^2/U^2}
       =\frac{V^2-z^2}{U^2-z^2}.
\]
Its coefficient \(V^2/U^2-1\) is positive, so
\(\widetilde\Phi>0\). The same decay permits arbitrary complex
heat time. Finally, expanding \(\cosh(x\sqrt w)\) shows that its
normalized coefficients are the strictly positive integrals
\[
 \frac{1}{(2j)!\widetilde H_{-a}(0)}
   \int_0^\infty x^{2j}e^{-a x^2}\widetilde\Phi(x)\,dx.
\]
\end{proof}

This proves preservation of the explicitly stated positive-kernel
heat class. The stronger definition of an admissible kernel in
\cite[Definition~1.2, p.~2]{Csordas2015} also requires strict decrease
of the transformed kernel; that additional property is not used here.
The exact Dirichlet-series obstruction and the transported positive
derivative-moment asymptotics remain those proved above for the same
rational replacement.

\begin{proposition}[The exact heat threshold after replacement]
\label{prop:fd-surgery-threshold}
The real times at which every zero of \(\widetilde H_t\) is real
are exactly \([0,\infty)\). Thus this transformed heat family has
Newman threshold zero and a double endpoint zero at \(V\).
\end{proposition}

\begin{proof}
We first verify forward preservation directly for the functions in
this family. Fix a real time \(r\) whose endpoint
\(f=\widetilde H_r\) has only real zeros. The positive kernel gives
\(f(0)>0\), and the preceding tail estimates imply
\[
 \log\max_{|z|\le R}|f(z)|=O_r(R\log(R+2)).
\]
For example, split the defining integral at a constant multiple of
\(\log(R+2)\); on the remaining tail the negative
\(-c e^{4x}\) term dominates \(Rx+|r|x^2\).
Evenness and the Hadamard factorization therefore give
\[
 \frac{f(z)}{f(0)}=\prod_{j\ge1}(1-z^2/x_j^2),
 \qquad \sum_jx_j^{-2}<\infty,
\]
where positive real zeros are listed with multiplicity.
Let \(p_M\) be a finite initial product. For every complex \(z\),
\[
 |p_M(z)|\le\prod_{j\ge1}(1+|z|^2/x_j^2)
             =\frac{f(i|z|)}{f(0)}.
\]

Every polynomial \(p+c p'\), with \(c\) real, has only real zeros
when \(p\) does. Indeed, away from the roots the equation is
\(p'/p=-1/c\); its left side is strictly decreasing between
successive distinct real roots and has the required infinite
endpoint limits. A root of multiplicity \(m\) retains
multiplicity \(m-1\), and counting the remaining roots proves
the assertion. The case \(c=0\) is immediate.
Consequently
\[
 (1-tD^2/n)^n p_M
 = (1-\sqrt{t/n}\,D)^n(1+\sqrt{t/n}\,D)^n p_M
\]
has only real zeros for \(t\ge0\).
Its coefficient limit as \(n\to\infty\) is
\(e^{-tD^2}p_M\), so this polynomial also has only real zeros.

For a standard real Gaussian \(Z\), the polynomial heat identity is
\[
 e^{-tD^2}p_M(z)=\mathbb E\,p_M(z+i\sqrt{2t}\,Z).
\]
The product bound and the growth estimate dominate the integrand
by a Gaussian-integrable function, uniformly for \(z\) in a compact
set and \(t\) in a bounded nonnegative interval. Passing to the
limit \(M\to\infty\) therefore gives
\[
 \lim_{M\to\infty}e^{-tD^2}p_M(z)
 =\frac{\mathbb E f(z+i\sqrt{2t}\,Z)}{f(0)}
 =\frac{\widetilde H_{r+t}(z)}{f(0)}.
\]
The last identity follows by integrating the positive-kernel
representation against the Gaussian; its expectation multiplies
\(\cos(zx)\) by \(e^{t x^2}\), and the same decay justifies Fubini.
The limit is nonzero, since \(\widetilde H_{r+t}(0)>0\).
Hurwitz's theorem on each half-plane consequently proves forward
preservation of real zeros. At \(r=0\), all zeros are real by the
exact endpoint replacement, so every \(t\ge0\) has that property.

At \(V\), the endpoint has multiplicity exactly two. Write
\(g=\widetilde H_0''(V)/2\ne0\).
Joint analyticity of the heat family and
\(\partial_t\widetilde H_t=-\partial_z^2\widetilde H_t\) give
\[
 \frac{\widetilde H_{-a}(V+\sqrt a\,w)}{a g}
          =w^2+2+O(\sqrt a),\qquad a\downarrow0,
\]
uniformly on compact \(w\)-sets.
Rouch\'e's theorem near \(w=\pm i\sqrt2\) gives nonreal zeros
for every sufficiently small negative time.
An earlier negative time with only real zeros would propagate
forward, by the argument just proved, to one of these small
negative times. This contradiction excludes every negative time.
\end{proof}

\subsubsection*{The exact change in the trace certificate}

The replacement also has an explicit action on the spectral and
arithmetic data. In the notation of
Theorem~\ref{thm:fd-arithmetic-complete}, its counted product moments
and distinct-coordinate trace obey
\[
 \widetilde S_k=S_k+2v^{-2k}-2u^{-2k},\qquad
 \operatorname{tr}(\widetilde C^{2k})
     =\operatorname{tr}(C^{2k})-2u^{-2k}.
\]
The first identity counts the removed pair and the added copy of
the pair at \(v\). In the second, the coordinate at \(v\) was
already present, while the coordinate at \(u\) disappears.
Subtracting gives the exact limiting-defect change
\[
 \widetilde D_{\infty,k}=D_{\infty,k}+2v^{-2k}.
\]
Here the two Hilbert spaces retain their own weighted coordinate
norms; this calculation compares the displayed traces directly.

On their common compactly supported test-function space, write
\(F=\widehat f\), \(G=\widehat g\). The corresponding Weil forms
satisfy
\[
 \widetilde q(f,g)-q(f,g)
 =\sum_{\varepsilon\in\{\pm1\}}
   \bigl(F(\varepsilon v)\overline{G(\varepsilon v)}
        -F(\varepsilon u)\overline{G(\varepsilon u)}\bigr).
\]
For \(h=f*\widetilde g\), Fourier multiplication therefore gives
the equivalent arithmetic correction
\[
 \widetilde W(h)-W(h)
   =2\int_{\mathbb R}h(x)
                \bigl(\cos(vx)-\cos(ux)\bigr)\,dx.
\]
Thus each finite Gram matrix changes by a sum of four explicitly
given evaluation tensors, of rank at most four. The correction
survives in the limiting multiplicity defect with the positive
value just computed. These identities specify the finite arithmetic
information changed by the surgery, alongside the positive-kernel
and leading-moment information that it preserves.


\subsection{The density quantifier for local finite tests}

For a positive simple endpoint $x$, choose a conjugation-invariant disk
containing no other zero of $H_0$.  Its unique zero remains real and
simple for a sufficiently short real time interval about zero.
The normalized polynomial of the corresponding squared-variable
contour has degree one, $1+\beta_x(a)w$, with $\beta_x(a)>0$.
Every one of its Toeplitz minors is nonnegative and every Jensen
polynomial is hyperbolic throughout that interval.  This follows
from its positive linear factor, or from the exact differential
identity $(1+\beta_x(a)D)z^N=z^{N-1}(z+N\beta_x(a))$.
Consequently every proved positive-density set of simple endpoints
provides the same-density family of local tests with infinite detecting
order on their respective persistence intervals.

A quantitative persistence theorem makes this uniform on the set of
labels to which it applies.  Specifically, a real simple zero branch
on $[-\tau,0]$ defines the same positive degree-one local factor for
all $a\in[0,\tau]$; this identity transfers the entire persistence
interval to the local detecting test.  The persistence theorem elsewhere
in this paper supplies its stated common value $\tau=c/\log^2T$.

The global quantities of Theorem~\ref{thm:fd-limits} use the product
over every endpoint.  Their full scaled limits give an exact quantifier
check: any sequence $a_n\downarrow0$ on which
$\sqrt{a_n}m_*(-a_n)\to\infty$, or
$a_nj_*(-a_n)\to\infty$, already proves simplicity of every endpoint.
The local degree-one factors enter that global product by multiplication;
a single multiple factor contributes the finite limiting constant
computed in Theorem~\ref{thm:fd-limits}.  This product map is the
precise relation between the positive-density local conclusion and
the all-endpoint growth problem.
